\section{Applicazioni Lineari (Nucleo e Immagine)}

\begin{concept}{Richiami di Teoria: Ker e Im}
Sia data l'applicazione lineare $f: V \to W$, determinata da una sua matrice associata $A$ tale che $L_A(x)=Ax$.
\begin{itemize}
    \item \textbf{Nucleo} ($\ker A$): L'insieme delle pre-immagini che restituiscono una trasformazione nulla. Le trovi ponendo $A \cdot x = 0$ e risolvendo tutto il sistema con Gauss.
    \item \textbf{Immagine} ($\text{Im } A$): L'area d'arrivo di tutti i vettori trasformati dallo span di $V$. È sempre generata dalle COLONNE di $A$ (colonne linearmente indipendenti formano la base per l'immagine).
\end{itemize}
\textbf{Teorema Rank-Nullity (Dimensione):} $\dim V = \dim(Ker(f)) + \dim(Im(f))$.
\end{concept}

\begin{logicbox}
\textbf{Pattern 4: Lo schema d'attacco classico per le Applicazioni}\\
Di fronte a una matrice che rappresenta l'applicazione lineare:
1. Per trovare l'immagine: Riduci la matrice con Gauss, guarda in quali colonne si creano i pivot. Quelli estratti dalle colonne corrispondenti ma della **matrice matrice originale di partenza** saranno i vettori della base dell'immagine!
2. Per trovare il Kernel: Usa la matrice precedentemente ridotta fino in fondo, ponila a zero ($Ax=0$) e risolvi a ritroso le equazioni in funzione delle variabili libere.
\end{logicbox}

\begin{exercisebox}[title={Esercizio 1: Determinare Ker e Im}]
Data la matrice $A$ dell'applicazione lineare:
$$ A = \begin{pmatrix} 1 & 1 & 0 \\ 0 & -2 & 1 \\ -1 & 1 & -1 \end{pmatrix} $$
Determinare una base per il Nucleo ($\ker A$) e una base per l'Immagine ($\text{Im } A$).
\end{exercisebox}

\begin{solbox}
\textbf{Studio della Matrice e Immagine:}\\
Riduciamo a scalini:
$$ A = \begin{pmatrix} 1 & 1 & 0 \\ 0 & -2 & 1 \\ -1 & 1 & -1 \end{pmatrix} \xrightarrow{R_3 \to R_3 + R_1} \begin{pmatrix} 1 & 1 & 0 \\ 0 & -2 & 1 \\ 0 & 2 & -1 \end{pmatrix} \xrightarrow{R_3 \to R_3 + R_2} \begin{pmatrix} 1 & 1 & 0 \\ 0 & -2 & 1 \\ 0 & 0 & 0 \end{pmatrix} $$
Il rango di $A$ è $r=2$. I pivot si trovano nelle colonne 1 e 2.
Per cui, le corrispondenti colonne della matrice originaria formano la base dell'immagine:
$$ \mathcal{B}_{Im} = \left\{ \begin{pmatrix} 1 \\ 0 \\ -1 \end{pmatrix}, \begin{pmatrix} 1 \\ -2 \\ 1 \end{pmatrix} \right\} \quad \implies \dim(Im(f)) = 2 $$

\textbf{Studio del Nucleo:}\\
Usiamo il teorema di nullità+rango: la dimensione di $\mathbb{R}^3$ del dominio \ è $3$. Allora $\dim(Ker) = 3 - 2 = 1$. Il nucleo dipende da $1$ parametro libero.
Risolviamo l'omogeno di base per lo zero usando la matrice a scalini derivata:
$$ \begin{cases} x + y = 0 \\ -2y + z = 0 \end{cases} $$
Sia $z = t$. Allora $y = t/2$. Di conseguenza $x = -t/2$.
$$ \ker(A) = \{ (-t/2, t/2, t)^T \} $$
Scegliamo arbitrariamente $t=2$ per eliminare le frazioni:
$$ \mathcal{B}_{Ker} = \left\{ \begin{pmatrix} -1 \\ 1 \\ 2 \end{pmatrix} \right\} $$
\end{solbox}
